%========================================
% [3] & [4] Pubbabhāganamakārapāṭha
%========================================
\section*{e3 \& e4 | Pubbabhāganamakārapāṭha}

\begin{chantsubsection}
  \chantnote
    {\{Handadāni mayaṃ taṃ bhagavantaṃ vācāya abhigāyituṃ pubbabhāganamakārañceva buddhānussatinayañca karomase\}}
  \chantnotetrans
    {\{Now, in order to praise the Blessed One with speech, we perform the preliminary homage and the method of recollecting the Buddha.\}}
\end{chantsubsection}
\chantgap
\begin{chantsubsection}
  \chantnote
    {Namo tassa bhagavato arahato sammāsambuddhassa}
  \chantnote{(three times)}
  \chantgap
  \chantnotetrans
    {Homage to the Blessed One, the Arahant, the Perfectly Enlightened One.}
\end{chantsubsection}
\chantgap
\begin{chantsubsection}
  \chantlineinline
    {Taṃ kho pana bhagavantaṃ evaṃ kalyāṇo kittisaddo abbhuggato:}
    {Now, about that Blessed One, this fine report has spread abroad:}
  \chantlineinline
    {Iti pi so bhagavā arahaṃ sammāsambuddho}
    {“Such indeed is that Blessed One: an arahant, perfectly awakened;}
  \chantlineinline
    {vijjācaraṇasampanno sugato lokavidū anuttaro}
    {accomplished in knowledge and conduct, well-gone, knower of the worlds, unsurpassed;}
  \chantlineinline
    {purisadammasārathi satthā devamanussānaṃ buddho bhagavā’ti}
    {trainer of persons to be tamed, teacher of devas and humans, the Buddha, the Blessed One.”}
\end{chantsubsection}
